\section*{Capítulo \ref{ch:b3:hydrogen-atom} --- El átomo de hidrógeno}

\begin{solution}{exo:b3:hydrogen-atom:1}
(a) $u'' = (1/a^2 - 2/ar)u$; los términos en $1/r$ coinciden cuando
$\hbar^2/m_{\text{e}}a = k$, es decir, $a = a_0$, y los términos constantes
dan $E = -\hbar^2/2m_{\text{e}}a_0^2 = -E_{\text{I}}$. (b) $C =
2/a_0^{3/2}$. (c) $[\hbar^2/mk] = \unit{m}$; $[mk^2/\hbar^2] =
\unit{J}$. (d) El argumento variacional del
\cref{exo:b3:schrodinger-three-dimensions:12} muestra que $-E_{\text{I}}$
es la energía mínima que puede alcanzar cualquier estado normalizado: el
principio de incertidumbre pone suelo al átomo.
\end{solution}

\begin{solution}{exo:b3:hydrogen-atom:2}
(a) \qty{121.6}{nm}; límite \qty{91.2}{nm}. (b) $656.3$ (roja),
$486.1$ (azul verdosa) y \qty{434.0}{nm} (violeta); límite
\qty{364.6}{nm} --- las tres primeras son visibles. (c) Desde
\qty{1875}{nm} hasta \qty{820}{nm}: infrarrojo cercano. (d)
$R_{\text{H}}c\,(1/109^2 - 1/110^2) = \qty{5.01}{GHz}$.
\end{solution}

\begin{solution}{exo:b3:hydrogen-atom:3}
(a) $\dd(r^2\eu^{-2r/a_0})/\dd r = 0$ en $r = a_0$. (b) $\langle r\rangle = (4/a_0^3)\int_0^\infty
r^3\eu^{-2r/a_0}\dd r = (4/a_0^3)\times 3!\,(a_0/2)^4 = \tfrac32
a_0$. (c) Con $x = 4$:
$\tfrac12(16 + 8 + 2)\eu^{-4} = 0.24$: una cuarta parte del tiempo el
electrón está más allá de $2a_0$. (d) El electrón no tiene radio: $P(r)$
es una distribución con moda, media y colas largas --- el átomo
es difuso al nivel de un factor dos.
\end{solution}

\begin{solution}{exo:b3:hydrogen-atom:4}
(a) $l = 0$: uno; $l = 1$: tres; $l = 2$: cinco --- nueve. (b) $2$,
$8$ y $18$. (c) Exactamente las longitudes de las tres primeras filas: la
tabla periódica es el registro del llenado de estas capas. (d) La
degeneración en $m$ sobrevive (el espacio sigue siendo isótropo); la
degeneración en $l$ se rompe porque el potencial apantallado que ve el
electrón de valencia ya no es $1/r$ puro ---
\cref{exo:b3:hydrogen-atom:8}.
\end{solution}

\begin{solution}{exo:b3:hydrogen-atom:5}
(a) Hidrogenoide con $Z - 1$: $h\nu = (Z-1)^2E_{\text{I}}(1 -
\tfrac14)$. (b) $784 \times \qty{10.2}{eV} = \qty{8.0}{keV}$ ---
coincide con la K$\alpha$ medida del cobre. (c) Que el entero que
ordena los elementos es la \emph{carga nuclear} y no el peso
atómico: los huecos en las líneas de Moseley localizaron elementos aún no
descubiertos. (d) $42^2 \times 10.2 = \qty{18.0}{keV}$ --- el tecnecio,
hallado décadas después, obedeció.
\end{solution}

\begin{solution}{exo:b3:hydrogen-atom:6}
(a) Dos masas iguales: $\mu = m_{\text{e}}/2$: energía de enlace de
\qty{6.8}{eV} y radio $2a_0 = \qty{0.106}{nm}$. (b) $\tfrac34
\times 6.8 = \qty{5.1}{eV}$: \qty{243}{nm}. (c) Dos fotones de
\qty{511}{keV}, en sentidos opuestos (conservación del momento en reposo).
(d) La tomografía por emisión de positrones: la pareja de fotones
colineales de \qty{511}{keV} es la señal que triangula todo anillo PET.
\end{solution}

\begin{solution}{exo:b3:hydrogen-atom:7}
(a) $a_{\text{ex}} = 0.0529 \times 12.9/0.058 = \qty{11.8}{nm}$.
(b) $E = 13.6 \times 0.058/12.9^2 = \qty{4.7}{meV}$. (c) El
tamaño natural del par supera al del punto: es la pared, y no la atracción,
la que da forma al estado --- el $1/R^2$ vence al $1/R$, como se prometió.
(d) A temperatura ambiente, $k_{\text{B}}T$ es cinco veces la energía de enlace: los
pares se ionizan; la física excitónica vive por debajo de
$\sim\qty{50}{K}$ en cristales limpios.
\end{solution}

\begin{solution}{exo:b3:hydrogen-atom:8}
(a) El $3s$: su lóbulo más interno se sumerge dentro del tronco, donde el
$+11$ nuclear apenas está apantallado. (b) $-13.6/9 = \qty{-1.51}{eV}$
coincide casi exactamente con el $3d$: el $3d$ orbita fuera del tronco y ve
un $+1$ neto; el $3s$ (y en parte el $3p$) prueban el potencial profundo y
se hunden. (c) $\Delta E = \qty{2.10}{eV}$: $\lambda = \qty{590}{nm}$
--- el amarillo del sodio de las farolas. (d) $n_{\text{ef}} =
\sqrt{13.6/|E|}$: $1.63$ y $2.11$, de modo que $\delta_s = 1.37$ y
$\delta_p = 0.89$.
\end{solution}

\begin{solution}{exo:b3:hydrogen-atom:9}
(a) Aproximadamente la energía de enlace de $n = 100$, $\sim\qty{1.4}{meV}$, más la
pequeña energía térmica del electrón: un fotón de infrarrojo lejano o de
radio. (b) Los pasos con $\Delta n = 1$ cerca de $n \approx 100$ caen en la
banda de radio de los GHz --- las líneas de recombinación. (c) H$\alpha$ es el paso $3 \to
2$, cerca del final de la cascada. (d) Lyman $\alpha$, aunque
es la más intensa, es ultravioleta y, en una nebulosa llena de hidrógeno en
el estado fundamental, se dispersa resonantemente y queda atrapada;
H$\alpha$ escapa sin trabas y pinta de rojo la nebulosa.
\end{solution}

\begin{solution}{exo:b3:hydrogen-atom:10}
(a) $\langle 1/r\rangle = 1/a_0$: $\langle E_p\rangle = -k/a_0 =
\qty{-27.2}{eV} = 2E_1$. (b) $\langle E_k\rangle = E_1 - \langle
E_p\rangle = +\qty{13.6}{eV}$: la razón del virial $-\tfrac12$ de
toda órbita en $1/r$. (c) $v = \sqrt{2E_k/m_{\text{e}}} = \alpha c =
\qty{2.2e6}{m/s}$. (d) $I = ev/2\pi a_0 \approx \qty{1}{mA}$
circulando a \qty{53}{pm}: $B \sim \mu_0I/2a_0 \approx
\qty{12}{T}$ --- el enorme campo interno del que se alimentará la
estructura hiperfina.
\end{solution}

\begin{solution}{exo:b3:hydrogen-atom:11}
(a) Del \cref{exo:b3:hydrogen-atom:10}(c), $v/c = \alpha$;
al escalar: $Z\alpha/n$. (b) $\alpha^2E_{\text{I}} = \qty{0.72}{meV}$;
en $n = 2$ los desdoblamientos salen de decenas de $\unit{\micro eV}$ ---
unos \qty{10}{GHz}, medibles por radio (y medidos). (c) $Z\alpha =
0.67$: dos tercios de la velocidad de la luz --- el tratamiento de
Schrödinger falla y toma el relevo la ecuación de Dirac. (d) Cuando
$Z\alpha \to 1$, la energía del estado fundamental se hunde formalmente por
debajo de $-2m_{\text{e}}c^2$: el propio vacío chisporrotearía pares
electrón--positrón --- campos supercríticos, buscados en las colisiones de
iones pesados.
\end{solution}

\begin{solution}{exo:b3:hydrogen-atom:12}
(a) Derivando: $\dot{\vect p}\wedge\vect L = -(mk/r^3)\vect
r\wedge(\vect r\wedge\dot{\vect r}\,m)$; al desarrollar el doble producto
vectorial se obtiene exactamente $mk\,\dd\vect e_r/\dd t$: $\dot{\vect A} =
\vect 0$. (b) Un eje conservado es una simetría más allá de la isotropía;
la degeneración es un etiquetado incompleto, y la etiqueta adicional (la
prima cuántica de $\vect
A$) conecta las $l$ dentro de un mismo $n$. (c) Con
$V = -k/r + \epsilon/r^2$, el momento angular efectivo se desplaza,
el periodo angular deja de coincidir con el radial y el eje de la
elipse gira con un ritmo $\propto\epsilon$. (d) Las correcciones
relativistas hacen el papel de $\epsilon$ en el propio hidrógeno: la
estructura fina del \cref{exo:b3:hydrogen-atom:11} es precisamente la
ruptura de la degeneración en $l$.
\end{solution}

\begin{solution}{pb:b3:hydrogen-atom:1}
\textbf{1.} $\langle r\rangle \sim n^2a_0$; $|E_n| =
E_{\text{I}}/n^2$; espaciado $\approx 2E_{\text{I}}/n^3$; frecuencia
orbital $\propto 1/n^3$ (Kepler sobre la órbita de Bohr).
\textbf{2.} $E_{n+1} - E_n \to 2E_{\text{I}}/n^3$, y $h$ por
la frecuencia clásica $v_n/2\pi r_n$ da la misma expresión
(el cálculo del \cref{pb:b3:hamiltonian-mechanics:1}, apartado 20).
\textbf{3.} $n = 50$: $\qty{132}{nm}$, \qty{5.4}{meV}; $n = 100$:
$\qty{0.53}{\micro m}$, \qty{1.4}{meV}.
\textbf{4.} $n^2ea_0 = \qty{2.1e-26}{C.m} \approx 6300\,\text{D}$:
el dipolo de tres mil moléculas de agua sobre un solo átomo.
\textbf{5.} $F \sim \num{5.4e-3}/\num{1.32e-7} \approx
\qty{4e4}{V/m} = \qty{400}{V/cm}$; el umbral exacto de ionización clásica
es unas ocho veces menor ($\sim\qty{50}{V/cm}$):
en cualquier caso, un susurro de campo.
\textbf{6.} Con secciones eficaces geométricas $\sim\pi n^4a_0^2$, incluso
el vacío ultraalto de laboratorio hace chocar a un átomo así en
microsegundos; las densidades interestelares ($\sim\num{e6}$ partículas por
$\unit{m^3}$) le dejan días --- el espacio es mejor cámara de vacío.
\textbf{7.} $\nu = R_{\text{H}}c\,(1/109^2 - 1/110^2) =
\qty{5.01}{GHz}$.
\textbf{8.} Radio centimétrica: un radiotelescopio --- una antena grande y
un receptor silencioso.
\textbf{9.} $\lambda = \qty{6.0}{cm}$; el átomo con $n = 110$ mide
$\sim\qty{0.6}{\micro m}$: cien mil átomos por longitud de
onda --- holgadamente un régimen de ``antena''.
\textbf{10.} $\Delta\nu/\nu \sim v/c = \num{4.3e-5}$: unos
\qty{200}{kHz} sobre \qty{5}{GHz} --- y la anchura medida devuelve
la temperatura de la nebulosa.
\textbf{11.} La Galaxia interior ionizada: regiones HII escondidas tras
polvo que extingue todo fotón de Balmer --- las líneas de recombinación de
radio cartografiaron la estructura espiral que la astronomía óptica no podía
ver.
\textbf{12.} Los microcampos eléctricos de los iones vecinos (ensanchamiento
Stark) difuminan los niveles, y la radiación de cuerpo negro y cósmica, más
las colisiones, ionizan o mezclan en $l$ esos estados frágiles.
\textbf{13.} $C_6 \sim d^4/\Delta E \propto (n^2)^4/n^{-3} =
n^{11}$: súbase $n$ de $1$ a $70$ y la interacción crece
veinte órdenes de magnitud.
\textbf{14.} Dentro de $R_{\text{b}}$, el estado de la pareja queda fuera
de resonancia, de modo que la excitación de un átomo \emph{condiciona} la
respuesta de su vecino: exactamente la lógica controlada que necesita una
puerta de dos cúbits.
\textbf{15.} Unas interacciones de alcance micrométrico permiten que cada
átomo ocupe su propia pinza direccionable, muy lejos según los cánones
atómicos y aun así fuertemente acoplado --- el alcance de la interacción
ajustado a la resolución óptica.
\textbf{16.} $\sim 200$ puertas por vida media: la razón acota la fidelidad
alcanzable, y es la razón, y no los microsegundos en bruto, de lo que vive
un procesador.
\textbf{17.} A \qty{300}{K}, el baño de fotones es denso en torno a
$\sim\qty{0.1}{meV}$: baraja e ioniza los niveles de Rydberg dentro de
sus vidas medias radiativas; unas pantallas frías (\qty{4}{K}) dejan sin
alimento esas transiciones.
\textbf{18.} Los ritmos de transición escalan como $d^2 \propto n^4$: en $n =
50$, seis millones de veces el acoplamiento de un átomo ordinario ---
razón por la cual una celda de vapor de átomos de Rydberg es hoy un sensor
calibrado de campos de microondas y de terahercios.
\textbf{19.} $92^2 \times \qty{13.6}{eV} = \qty{115}{keV}$ hasta
$13.6/10^6\,\unit{eV} = \qty{13.6}{\micro eV}$: diez órdenes de
magnitud a partir de una sola fórmula.
\textbf{20.} La clasicidad amanece donde el espaciado se convierte en una
fracción evanescente de la energía, $n \gg 1$ --- el límite de
correspondencia propio del átomo.
\textbf{21.} Es el único átomo calculable desde primeros principios
hasta la precisión del experimento: toda discrepancia es un descubrimiento,
de modo que el hidrógeno ancla las constantes fundamentales.
\textbf{22.} $\tau \approx 1.6\,\unit{ns} \times 50^3 =
\qty{0.2}{ms}$ --- generosa, pero la redistribución por cuerpo negro (parte
III) se la come, otra razón para las pantallas frías.
\textbf{23.} La simetría CPT --- los átomos de materia y de antimateria
deben coincidir línea a línea; el hidrógeno es el ancla porque solo ahí
llega la teoría a la decimoquinta cifra junto con el experimento.
\textbf{24.} El dipolo $n^2$ hace que el átomo sienta el campo de un solo
fotón de microondas, mientras que su larga vida media y su estructura de
niveles le permiten adquirir una fase medible \emph{sin} absorber el fotón:
detección cuántica de no demolición.
\textbf{25.} Las dependencias $n^2$, $n^{-2}$, $n^{-3}$ y $n^{11}$ estiran
el único átomo resuelto de \qty{52.9}{pm} a medio micrómetro y
de \qty{121.6}{nm} a \qty{6}{cm} --- la misma solución de Schrödinger
emitiendo desde las regiones HII de la Galaxia y marcando el ritmo de las
puertas de dos cúbits en arreglos de pinzas ópticas.
\end{solution}
